\cotefile{27}
\scanpage{2}
Divers pour\\
$\overline{\text{\uline{XII}}}$ à $\overline{\text{\uline{XIV}}}$
% archive sheet number: 1
\scanpage{3}
\mspage{\handwritten{6}}
\margin{\handwritten{Co Exp $\overline{\text{\uline{XIII}}}$}}

\uline{Proposition} 1.3. Avec les notations précédentes, \deleted{\unsure{en|on}} supp\handwritten{o}s\deleted{\unsure{ant|e}}\handwritten{ons} \deleted{encore} \deleted{\unsure{la|la en}} l'inégalité (2) valable\deleted{,}\inserted{et $G,V,W$ géométriquement irréductibles.} \deleted{c}\handwritten{C}onsidérons les conditions suivantes:

(i) $\handwritten{N}=\handwritten{M}_{\deleted{\unsure{a}}a}$, et $N$ est lisse sur $k$.

(ii) $\handwritten{N}=\handwritten{M}_{a}$, et $\dim X=\dim V$ (cf lemme 1.2.).

(iii) $\handwritten{N}=\handwritten{M}_{a}$, et $\handwritten{\psi}:X\handwritten{\to}V$ est génériquement quasi-fini.

(iii bis) $\handwritten{N}=\handwritten{M}_{a}$, et $\handwritten{\psi}:X\handwritten{\to}V$ est génériquement non ramffié.

(iii ter) $\handwritten{N}=\handwritten{M}_{a}$, et $\handwritten{\psi}:X\handwritten{\to}V$ est génériquement étale.

(iv) $\handwritten{\psi}:X\handwritten{\to}V$ est non ramffié en $(e^{\handwritten{\unsure{\circ|\bullet}}},a)$.

(iv bis) $\handwritten{\psi}:X\handwritten{\to}V$ est étale en $(e^{\handwritten{\circ}},a)$.

(v) $\handwritten{\psi}:X\handwritten{\to}V$ est lisse en $(e^{\handwritten{\circ}},a)$, et $N$ est lisse sur $k$, (2) une \handwritten{égalité.} \deleted{égalité}

(vi) $M_a$ et $N$ coïncident au voisinage de $e$.

Alors on a les implications
\[
\resizebox{.88\linewidth}{!}{$\displaystyle
\text{(\deleted{\unsure{i}}v)}\handwritten{\Leftrightarrow}\text{(iv bis)}\handwritten{\Longrightarrow}
\text{(iii \deleted{bis}ter)}\handwritten{\Longrightarrow}\text{(i)}\handwritten{\Leftrightarrow}
\text{(ii)}\handwritten{\Leftrightarrow}\text{(iii)}
\overset{\inserted{\handwritten{\Leftrightarrow\text{(iii bis)}\Leftrightarrow}}}{\mathstrut}
\text{(iv)}\handwritten{\Leftrightarrow}\text{(vi)}\handwritten{\ ;}$}
\tag*{\handwritten{(3)}}
\]
si \deleted{\unsure{$G$|$V$}}\handwritten{$V$} est \deleted{lisse sur $k$ et $W$ géométriquement} \handwritten{\inserted{génériquement}} réduit, on a (iii \deleted{bis}\handwritten{ter})$\handwritten{\Leftrightarrow}\text{(\handwritten{i})}$\handwritten{,} si \deleted{de plus}\handwritten{\inserted{enfin}} $V$ est \deleted{géométriquement} normal \deleted{\illegible{1}}\handwritten{\inserted{on a}} toutes les conditions envisagées sont équivalentes.

On sait déjà par le lemme 1.1. que l'on a $\text{(i)}\handwritten{\Rightarrow}\text{(iv)}\handwritten{\Leftrightarrow}\text{(v\handwritten{i})}$, d'autre part (vi) implique évidemment par définition $\handwritten{N}=\handwritten{M}_a$ et (iv) implique \deleted{(iii bis)} que $\handwritten{\psi}$ est génériquement non ramifié (NB il résulte des hypothèses faites que $X$ est irréductible\deleted{\unsure{s|,}} comme image de $G\mathrm{x}W$ qui l'est), donc $\text{(iv)}\handwritten{\Leftrightarrow}\text{(vi)}\handwritten{\Rightarrow}\text{(iii bis)}$, évidemment $\text{(iii bis)}\handwritten{\Rightarrow}\text{(iii)}$\unsure{,} \deleted{enfin} De plus $\text{(iii)}\handwritten{\Rightarrow}\text{(ii)}$, car comme $\dim X\handwritten{\geq}\dim V$ (1.2.) il s'ensuit que si $\handwritten{\psi}$ est génériquement quasi-fini, \deleted{\unsure{on a}}\handwritten{donc} $\dim X\handwritten{\leq}\dim V$\handwritten{,} \deleted{\unsure{donc}}\handwritten{on a} $\dim X=\dim V$. Enfin, en vertu de 1.2.\handwritten{,} $\text{(ii)}\handwritten{\Rightarrow}\text{(i)}$, ce qui prouve l'équivalence des conditions (i) (ii) (iii) (iii bis) (iv) (vi). Evidemment $\text{(iii ter)}\handwritten{\Rightarrow}\text{(iii bis)}$ et $\text{(iv bis)}\handwritten{\Rightarrow}\text{(iv)}$, donc par ce qui précède
% archive sheet number: 2
\scanpage{4}
\mspage{\handwritten{7}}
(iv bis) implique $\handwritten{N}=\handwritten{M}_a$, d'où aussitôt que $\text{(iv bis)}\handwritten{\Longrightarrow}\text{(iii ter)}$. Enfin, (iv bis) signifie aussi que $\handwritten{\psi}:X\handwritten{\to}V$ est lisse en $(e^{\handwritten{\bullet}},a)$ et que $\dim X=\dim V$, donc équivaut à (v) en vertu de 1.2. Cela prouve les implications (3). Si $V$ est \deleted{réduit (ou seulement} génériquement réduit\deleted{)}, \deleted{alors}\inserted{comme} (iii bis) implique que la fibré générique de $X\handwritten{\to}V$ est non ramifiée sur le corps rédiduel du point générique de $V$, donc étale sur ce dernier, et que ce corps est aussi l'anneau local de $V$ en son point générique, il s'ensuit que $X\handwritten{\to}V$ est génériquement étale, \deleted{\unsure{il|i.e.|et}} par suite on a \handwritten{(}iii bis)$\handwritten{\Rightarrow}\text{(iii ter)}$. Enfin, si $V$ est normal en $a$, alors $\text{(iv)}\handwritten{\Longrightarrow}\text{(iv bis)}$, car comme $\dim X\handwritten{=}\dim V$ par 1.2. , (iv) implique que $\handwritten{\psi}$ est dominant, et on peut appliquer SGA I 9.5. (ii) qui implique que $\handwritten{\psi}$ est étale et non seulement \inserted{non} ramifié en $(e^{\handwritten{\bullet}},a)$, \deleted{\unsure{ce qui prouv|ce qui prouve}} Cela achève la démonstration de 1.3. Notons d'ailleurs que la condition (ii) implique déjà en vertu de 1.2. que l'inégalité (2) est en fait une égalité. Notons aussi le

\uline{Corolla\deleted{\unsure{o}}ire} 1.4. Supposons que l'on ait $\handwritten{N}=\handwritten{M}_a$\handwritten{,} \inserted{$V$ normal en $a$.} Alors $\phi$ est lisse en $(e,a)$, \deleted{\unsure{(}}donc (1.1. b) $M_a$ est lisse sur $k$ en $e$, et $\deleted{\phi}\handwritten{\psi}$ est lisse en $(e^{\handwritten{\bullet}},a)$.

En effet, on a les inégalités
\[
\dim V+\deleted{\unsure{\dim|\mathrm{r}}}\operatorname{rang}\handwritten{N}\handwritten{\leq}
\dim(G\mathrm{x}W)\handwritten{\leq}\dim_{(e,a)}\phi^{-1}(a)+\dim V
\handwritten{\leq}\dim\handwritten{M}_a+\dim\deleted{\unsure{W|V}}\handwritten{V},
\]
où la première inégalité n'est autre que (2), la deuxième inégalité provient de la \deleted{\unsure{théorie}} théorie de la dimension pour le morphisme \deleted{\unsure{$\phi$}} $\phi:G\mathrm{x}W\handwritten{\to}V$ en le point $(e,a)$, et est une égalité si et seulement si $\phi$ est dominant et la dimension de $\phi^{-1}(a)$ en $(e,a)$ est celle de la fibre générique, ce qui (en vertu de l'hypothèse que $V$ est normal en $a$, et d'un théorème de Chevalley ) signifie aussi que $\phi$ est universellement ouvert en $(e,a)$ ; enfin la troisième inégalité \deleted{est} pro\deleted{\unsure{uv}}\handwritten{vie}nt de l'isomor%
% archive sheet number: 3
\scanpage{5}%
\mspage{\handwritten{8}}%
phisme $\phi^{-1}(a)\handwritten{\simeq}M_a$, transformant $(e,a)$ en $e$, \margin{\handwritten{e}}c'est une égalité si et seulement si $M_a$ est \inserted{\handwritten{régulier en \deleted{\unsure{$a$}}$e$ i.e.}} lisse sur $k$ en $e$ i.e. $\phi^{-1}(a)$ est lisse sur $k$ en $(e,a)$. Si \deleted{\unsure{donc}} $\handwritten{N}=\handwritten{M}_a$, alors les membres extrêmes de la suite d'inégalités précédente sont égaux, donc tous les termes sont égaux, ce qui implique que $\phi$ est universellement ouvert en $(e,a)$ et que sa fibre en ce point est lisse sur $k$. Comme $V$ est \deleted{réduit \unsure{en}} normal donc réduit en $e$, il \deleted{résulte \unsure{d'un}}\inserted{s'ensuit, (\handwritten{p}ar un} théorème de EGA IV) que $\phi$ est lisse \deleted{\unsure{sur}} en $(e,a)$, ce qui prouve 1.4. La démonstration invoque certains résultats assez délicats de \inserted{la} théorie des schémas, \deleted{\unsure{mais}} mais nous n'au\handwritten{\unsure{t|r}}ons pas besoin de 1.4. pour la suite.

Les conditions de 1.3. étaient relatives au choix d'un point $a$ de $\mathbb{W}(k)$. On conclut de 1.3. l'énoncé suivant, indépendant d'un tel choix :

\uline{Corollaire} 1.5. Soient $G,V,W$ sur $k$ comme dans 1.3. \deleted{Considérons les} et supposons $V$ génériquement \deleted{réduit. Alors les seul}\inserted{\handwritten{séparable sur $k$.}} Considérons les conditions suivantes:

(i) \deleted{\unsure{$\phi$}\illegible{1}} $\handwritten{\psi}:X\handwritten{\to}V$ est génériquement non ram\deleted{\unsure{ff|if}}\handwritten{if}ié.

(i bis) $\handwritten{\psi}:X\handwritten{\to}V$ est génériquement étale.

(i\deleted{\unsure{i}} \deleted{bis}\handwritten{ter}) $\handwritten{\psi}:X\handwritten{\to}V$ est génériquement lisse, $N$ est lisse sur $k$ et l'inégalité (2) une égalité .

(ii) \deleted{Il existe $a$ \unsure{$\in$} $W(k)$ tel que} $N$ est lisse sur $k$, et il existe $a\handwritten{\in}W(k)$ tel que $\handwritten{M}=\handwritten{N}_a$.

(iii) $\handwritten{\psi}:X\handwritten{\to}V$ est génériquement quasi-fini, et il existe $a\handwritten{\in}W(k)$ tel que $\handwritten{M}=\handwritten{N}_a$.

(iv) Il existe $a\handwritten{\in}W(k)$ tel que $M_a$ et $N$ coïncident au voisinage de $e$\handwritten{.}

Alors on a les implications
\[
\text{(ii)}\handwritten{\Leftrightarrow}\text{(iii)}\handwritten{\Leftrightarrow}\text{(iv)}
\handwritten{\Longrightarrow}\text{(i)}\handwritten{\Leftrightarrow}\text{(i bis)}\handwritten{\Leftrightarrow}\text{(i ter)}\ ,
\]
et si $k$ est algébriquement clos, toutes les conditions envisagées sont
% archive sheet number: 4
\scanpage{6}
\mspage{\handwritten{9}}
équivalentes.

L'équivalence de (ii) (iii) (iv) (en fait, pour tout $a\handwritten{\in}W(k)$ choisi) et le fait que ces conditions impliquent toutes les autres, \deleted{\unsure{e}}est contenu dans 1.3. En vertu de 1.2. le fait que $N$ est lisse sur $k$ et (2) une égalité équivaut à $\dim X=\dim V$, ce qui montre l'équivalence de (i bis) et (i ter), qui impliquent évidemment (i), reste à prouver que (i) implique (i bis), et , lorsque $k$ est algébriquement clos, que (i) implique (ii). Quitte à passer à la clôture algébrique de $k$, on peut supposer $k$ algébriquement clos , et on est réduit à prouver alors $\text{(i)}\handwritten{\Rightarrow}\text{(ii)}$. Or comme $\dim X\handwritten{\geq}\dim V$ en vertu de 1.2. , on voit que (i) implique que $\deleted{\phi}\handwritten{\psi}:X\handwritten{\to}V$ est dominant , et comme $V$ est génériquement réduit, \deleted{\unsure{il est génériquement régulier, il existe donc un ouvert}}\inserted{il e\handwritten{x}iste un ouvert non vide de $X$ qui est régulier et à} \unsure{fortiori|fortiooi} normal, donc il existe un \deleted{\unsure{$x\ X$ tel}} $x\handwritten{\in}X(k)$ tel que $\deleted{\phi}\handwritten{\psi}$ soit non ramifié en $x$ et que $\deleted{\phi}\handwritten{\psi}(x)$ soit un point \deleted{normal de $V$, \unsure{alors|Alors}} régulier et à fortiori normal de $V$. On aura $x=\deleted{\unsure{a|\phi}}(g^{\handwritten{\unsure{\circ|\bullet}}},a)$, et quitte à remplacer $x$ par $(e^{\handwritten{\circ}},g^{-1}.a)$, on peut supposer $a\handwritten{\in}W(k)$. Alors en vertu de(1.3. $\text{(iv)}\handwritten{\Rightarrow}\text{(i)}$ ) il s'ensuit que $a$ satifait à la condition 1.5. (ii), ce qui achève la démonstration.

\uline{Corollaire} 1.6. Soient $G,V,W$ comme dans 1.\deleted{\unsure{5}}\handwritten{3}. \deleted{\unsure{Il}}\inserted{\handwritten{$V$ normal, $k$} alg. clos.} Pour que $\handwritten{\psi}:X\handwritten{\to}V$ soit un morphisme \uline{birationnel}, il faut et il suffit que les conditions équivalentes de 1.5. soient satisfaites, et que de plus il existe un ouvert non vide $U$ de $V$ tel que tout $a\handwritten{\in}U(k)$ soit contenu dans un unique conjugué de $W$ par un élé\handwritten{ment} de $G(k)$. De plus, pour un $a\handwritten{\in}V(k)$, les conditions suivantes sont alors équivalentes:

(i) Il existe un unique conjugué de $W$ contenant $a$.

(ii) L'ensemble des conjugués de $W$ contenant $a$ est fini non vide.

(iii) Le schéma $\handwritten{\psi}^{-1}(a)$ \handwritten{``}des conjugués de $W$ qui contiennent $a$\handwritten{''} contient un point isolé.
% archive sheet number: 5
\scanpage{7}
\mspage{\handwritten{10}}
Lorsque \deleted{\unsure{$X$ et $W$ sont}} de plus \inserted{$W$ est} réduit\deleted{\unsure{s}}, et si $U$ désigne le plus grand ouvert de $V$ tel que $\deleted{\phi}\handwritten{\psi}$ induise un isomorphisme $\deleted{\phi}\handwritten{\psi}^{-1}(U)\handwritten{\to}U$, alors les conditions précédentes équivalent aussi à

(iv) $a\handwritten{\in}U(k)$.

La première assertion est claire sur la forme 1.5. (i bis).

L'équivalence de (i) à (iv) est une conséquence bien connue du \uline{Main Theorem} de Zariski\handwritten{.} \handwritten{On} not\handwritten{era} que \deleted{\unsure{la}} l'hypothèse que $\handwritten{\psi}$ est birationnel implique que \deleted{\unsure{$G/$ $X$, donc $G/N$, $X$ est,}} $X$ est génériquement réduit, donc (comme $p:X\handwritten{\to}G/N$ est fidèlement plat) que $G/N$ est génériquement réduit donc réduit\handwritten{;} par suite si $W$ est réduit, il en \deleted{\unsure{e}}est de même de $X$ qui est \deleted{fibré sur $G$} plat sur $G/N$ à fibres \unsure{géométriques|géométriqques} isomorphes\deleted{\unsure{à}} à $W$ donc réduites. En tous cas, que \handwritten{$W$} s\handwritten{oi}t réduit ou non, $X$ est irréduct\handwritten{i}ble (comme image de $\deleted{\unsure{X|G}}G\mathrm{x}W$ qui l'est, $G$ et $W$ l'étant), et l'équivalence de (i) à (iii) résulte alors du Main T\deleted{\unsure{he}}heorem appliqué au morphisme $X_{\text{réd}}\handwritten{\to}V$ induit par $\deleted{\phi}\handwritten{\psi}$.

\uline{Remarque} 1.7. Sous les conditions de 1.6. , un point $a$ de $V(k)$ qui satisfait aux conditions équivalentes (i) à (iii) pourra être appelé un \uline{point} \deleted{\unsure{W-rgéu}} $W$-\uline{régulier} de $V$ (relatif aux opérations de $G$ sur $V$). \deleted{Cette notion manifestement est stable par extension du corps de base $k\quad k'$, avec $k'$ algébriquement clos.} Si on ne suppose plus $k$ algébriquement clos, on dira encore qu\deleted{e} 'un point $a$ de $V$ est \inserted{$W$-}régulier si $\phi^{-1}(a)$ est \inserted{non vide et} fini sur $k(a)$\handwritten{.} \deleted{\handwritten{\unsure{i.e.}\ \illegible{3}}} Sous les conditions plus générales de 1.5. (i) à (i ter) , i.e. $\handwritten{\psi}$ étant génériquement étale (mais pas nécessairement birationnel), \deleted{\unsure{on}}il y a lieu d'introduire le degré de projection $d$ de $\handwritten{\psi}$, et le plus grand ouvert $U$ de $V$ tel que $\deleted{\phi}\handwritten{\psi}$ induise un morphisme $\deleted{\phi}\handwritten{\psi}^{-1}(U)\handwritten{\to}U$ qui soit \uline{étale} et \uline{fini}. Lorsque $V$ est normal\inserted{\handwritten{et} $X$ réduit (p.ex. $W$ réduit)}\handwritten{,} alors un $a\handwritten{\in}V$ est dans $U$ si et seulement si $\deleted{\phi}\handwritten{\psi}^{-1}(a)$ est fini et a un degré séparable égal à $d$ , i.e. (lorsque $k$ est algébriquement clos et $a$ est ra\handwritten{t}io\handwritten{nn}el
% archive sheet number: 6
\scanpage{9}
\uline{Compléments} $\overline{\text{\uline{XII}}}$

6. \uline{Le tore central maximal}. \uline{Relations avec le plus grand Tore quotient d'un \unsure{gpe|groupe} algébrique}

6.1. \uline{Définition en général} \unsure{Compatible|Compatibilité} aux extensions de la base descente, \unsure{générisation|spécialisation}\ldots\ etc\ldots

6.2. \uline{Existence lorsqu'il existe un \unsure{centr.|cart.|autre} réductif}

6.3. \deleted{\uline{N.B.} \unsure{En particulier, les groupes algébriques|En particulier, des groupes algébriques}} \uline{Un groupe algébrique : un \deleted{\unsure{tore}} maximal central} [N.B. le groupe algébrique n'a pas à être affine\ldots].

\deleted{\uline{6.4.} Soit $G$ un groupe algébrique \inserted{lisse \deleted{\unsure{connexe}}\unsure{connexe} affine} sur un corps parfait, \deleted{alors} soit $R$ son tore central maximal. Alors il existe un sous-groupe fini $\unsure{Z|3}$ de $R$, tel que \unsure{!} $\deleted{\unsure{R}}G/Z$ soit \deleted{facteur direct du} isomorphe au produit de $\deleted{G}R/Z$ et d'un groupe $H$. De plus, $H$ est tel que}

\margin{(\unsure{Rosenlicht})}
\uline{Th. 6.4} $X,Y$ \uline{géom. intègres} sur un corps $k$, munis de pts $a,b$ à valeurs dans $k$. Soit $\varphi:X\times Y\unsure{\mathord{-}|\to}G_m$ tel que $\deleted{\unsure{f}}\varphi(x,y)=1$, alors $\varphi$ est de la forme
\[
\varphi(x,y)=\varphi_1(x)\varphi_2(y)
\]
où $\deleted{\varphi_1(\varphi_2)}$ $\unsure{v|u}$ $(w)$ est un morphisme pointé $X\mathrel{-}G_m$ $(Y\mathrel{-}G_m)$.
% archive sheet number: 8
\scanpage{10}
\deleted{On peut supposer que \inserted{$k$ alg. clos} $X,Y$ affines.}
\margin{\unsure{!}}

Equivaut à ceci : \deleted{\unsure{de|$k$}} Si $u(x,\deleted{\unsure{a}}b)=1$, $u(\unsure{-|a},y)=1$ \unsure{identiqut}, alors $u(x,y)=1$ \unsure{identiqut}. Sous cette forme, on peut supposer \deleted{\unsure{à|$k$}} $k$ alg. clos.

\deleted{Alors \inserted{\unsure{Rien n'est|Bien entendu} \unsure{changé}} l'énoncé \unsure{précédent|équivalent} \unsure{consiste à|aussi à} \unsure{Notons que l'énoncé|N.B. que l'énoncé} \unsure{à remplacer|implique} $X,Y$ par Alors l'énoncé \unsure{vient de cet énoncé} : \unsure{équivaut|génériquement} à : \unsure{tout} \inserted{$\varphi(x,y)$} \unsure{hom} $X\times Y\unsure{\mathord{-}|\to}G_m$ \unsure{a la forme|est de la forme}}

On peut supposer $X,Y$ affines, et une réduction immédiate nous ramène au cas où $X,Y$ de dim 1.

\uline{Corollaire} Kif-kif, $G_m$ remplacé par un \deleted{\unsure{tore|des|deux}} \unsure{gpe|groupe} de \unsure{t. m.}

\deleted{$X\times Y\subset\overline X\times\overline Y$}
\deleted{$\unsure{\mathcal O|O}(\overline X\times\overline Y)$}

\uline{Corollaire} 1 \inserted{Soit $T$ un tore.} Tout morphisme \deleted{d'un} $u:G\mathrel{-}G_m$, \deleted{qui} $G$ \unsure{gpe|groupe} \uline{lisse} \uline{connexe}, qui envoie $e$ dans $e$, est un hom de groupes. \deleted{\unsure{De plus|D'ailleurs}}

\deleted{\unsure{Soit}} \deleted{\uline{Corollaire} 3} \deleted{\uline{Proposition} \unsure{2|3|?} $\mathrm{Hom}(G,G_m)$ est \unsure{un groupe de type fini.}} \deleted{\unsure{Prop}} $\mathrm{Hom}_k(G,G_m)$ \deleted{Déf} est un gpe de type fini.
% archive sheet number: 9
\scanpage{11}
\uline{Proposition} a) \inserted{$G$ \deleted{\unsure{disons}} \unsure{groupe} algébrique} Il existe un \deleted{\unsure{hom|hom.}}\inserted{hom} [\uline{épimorphisme}]
\[
G\longrightarrowH
\]
$H$ de t. m., tel que \deleted{\unsure{tt|$H$} \unsure{autre}} hom de $G$ dans un gpe de t. m. se factorise de façon unique par $H$.
% in the left margin
\diagram{\begin{tikzcd}[ampersand replacement=\&,row sep=0.6ex,column sep=0.5em]
 R \& \subset \& Z \& \deleted{\unsure{G}}\arrow[dd] \\
 \cap \& \& \cap \& {} \\
 T \& \subset \& G\arrow[d] \& {}\\
 {} \& {} \& \Gamma \& {}
\end{tikzcd}}

b) Lorsque \uline{$k$ est parfait}, la formation de $H$ est compatible aux \deleted{\unsure{tt|$H$}} extensions de la base.
\margin{$G$}

c) Lorsque $G$ est connexe \inserted{et lisse}, $H$ est un tore.
% in the left margin beside N.B.
\diagram{\begin{tikzcd}[ampersand replacement=\&,row sep=1em,column sep=1.7em]
 G\arrow[r,"\text{épim}"]\arrow[d,dash]\arrow[dr,phantom,"\deleted{\unsure{a}}" description] \& H\\
 G^{\circ}_{\text{réd}}\arrow[r,dash] \& H^{\circ}\arrow[u]
\end{tikzcd}}
\[
G/G^{\circ}_{\text{réd}}\unsure{\mathord{-}|\to}H/H^{\circ}
\]
\uline{N.B.} Si $k$ est imparfait, $\unsure{=|\simeq}$ si $G$ est commutatif, \inserted{\unsure{lisse sur le|linéaire et}\ \unsure{même|corps},} la formation de $H$ n'est en général pas compatible avec

\deleted{\uline{Corollaire}}\inserted{Proposition} $G$ un groupe\deleted{\unsure{$k$}} connexe \deleted{lisse} sur $k$, $R$ un \deleted{gpe de \unsure{t. b.|t. m.}}\inserted{tore}. Alors tte extension de $G$ par un $R$ qui splitte en tant que fibré principal homogène \uline{splitte}.
\[
\deleted{H^{2}_{\text{Hochschild}}(G,R)\simeq}
\]
\[
\left.\begin{aligned}
0\to C^{\bullet}(e,R)\to C^{\bullet}(H,R)&\longrightarrow
\deleted{E}\operatorname{Hom}_{\mathrm{gr}}(H^{\bullet},R)\mathrel{-}0\\
&H^i(\operatorname{Hom}_{\mathrm{gr}}(H^{\bullet},R))
\end{aligned}\right]
\]
\margin{En effet \inserted{\unsure{l'ext. est centrale|d'ext. centrale}} \unsure{Conjuguant|Composant} les deux résultats qui précèdent, elle provient d'une extension de $H$ par $R$, donc splitte.}
% archive sheet number: 10
\scanpage{12}
\margin{$\overline{\text{\uline{XII}}}$}
\uline{Cor} 5.8. Supposons $S$ \uline{normal} \deleted{. Soit} \inserted{\uline{noeth}.} irréductible. Soit $U$ un ouvert de $S$, \unsure{et} tel que $S-U$ de codim $\geq2$. Alors tt \deleted{tore sur} ss-groupe de t. m. de $G\mid U$ provient d'un unique sous-gpe de t. m. de $G$.

\uline{Cor.} 5.9. Supposons $S$ normal noeth irréd. de pt gén $g$, soit $\deleted{\unsure{H|\overline H}}T_g$ un \deleted{sous-groupe de t. m.}\inserted{sous-tore} de $G_g$. Pour que $H_g$ soit la fibre gén. d'un sous-\deleted{gpe de \unsure{t. m.}}\inserted{tore} $\deleted{\unsure{H|\overline H}}T$ de $G$, il f. et s. que pour tt $s\in S$ tel que $\deleted{\dim(\unsure{H|T|\overline H}_g)}$ $\dim\unsure{\mathcal O|\mathscr O}_{S,s}=1$, $\left(\overline{H_g}\right)^{\circ}_{s,\deleted{\unsure{\mathrm{réd}}}\mathrm{réd}}$ soit \deleted{un tore.} de t. m.
% slanted
\margin{\uline{rédiger} \uline{mieux}}

\deleted{\uline{Cor.} 5.10. Sous les conditions précédentes, supposons que $H_g$ soit $=$ un \uline{tore} [plus gén., que \deleted{\unsure{son type n'ait pas de}}\inserted{les \unsure{diviseurs} de la} \unsure{p-torsion|torsion} de son type soient $\unsure{\ne|\pm}$ car. résiduelles] \uline{Alors}}

On utilise un

\uline{Lemme} \ldots

\margin{TSVP}
% archive sheet number: 11
\scanpage{13}
On utilise un corollaire nouveau \unsure{?|:}
\margin{\unsure{!!}}

\margin{$\overline{\text{\uline{XI}}}$ \deleted{\unsure{XII|4}}}
\uline{Corollaire} 4.\deleted{\unsure{8}}\unsure{7|8}. \deleted{\uline{\unsure{base locale}}}

$G$ lisse affine sur $S$ local de pt fermé $s$. $\deleted{\unsure{H}\ \unsure{(H}}$ $H(n)$ famille de ss-gpes de $G$, de t. m., $H(n)$, $={}_nH(m)$, on suppose qu'il existe un ss-gpe de t. m. $H_s$ de $G_s$ \unsure{tels|tel} que $H(n)_s={}_nH_s$ pour tt $n$. Alors $\exists$ ss-gpe $H$ \inserted{de t. m.} unique de $G$ tel que $ {}_nH\unsure{=|\mathrel{:=}} \unsure{H_n|H(n)}$, $\forall n$.

\uline{Remarque} \deleted{Il suffit} Enoncé analogue en se bornant, aux $n$ dont l'un des diviseurs premiers est $\in P\ne\varnothing$ ($P$ ens de \unsure{nbres|nbs} premiers) pourvu que la torsion du type $H_s$ soit première à $P$. On utilise pour ceci la variante de la démonstration précédente, avec $M$ \unsure{courant|parcourant|remplacé} un \deleted{\unsure{sys}} ss-groupe de t. m. de $=$ type que $H_s$ [et en \deleted{\unsure{se bornant}} utilisant le th. de densité dans la dém. sous la forme \ldots]
% archive sheet number: 12
\scanpage{14}
\uline{Proposition} Soit $G$ \unsure{gpe|groupe} algébrique affine sur corps alg. clos $k$. Alors $\operatorname{Pic}(G)$ est fini. Si $G_{\mathrm{réd}}$ est \deleted{lisse et} \uline{résoluble}, \deleted{$\operatorname{Pic}(G)$ \unsure{est}} $\operatorname{Pic}(G)$ est nul. Si $G$ est lisse, \deleted{\illegible{1}}\inserted{connexe}, \deleted{$G$ est un $\operatorname{Pic}(G)$ est nul \unsure{ssi}} 
\[
\operatorname{Pic}(G)\simeq\operatorname{Pic}(G/B)/\operatorname{Im}D(T)
\]

\diagram{\begin{tikzcd}[ampersand replacement=\&,column sep=small,row sep=tiny]
R \& \subset \& B \& \subset \& G\arrow[r] \& G/B\\[-2pt]
R\cap T \& \& \overset{\cup}{T} \& \& \&
\end{tikzcd}}

[$B$ Borel, $T\subset B$ tore maximal,
\[
\underset{\substack{\scriptscriptstyle\parallel\\\operatorname{Hom}(B,G_m)}}{D(T)}
 \longrightarrow \operatorname{Pic}(G/B) \quad\text{hom caractéristique}\,]
\]

\diagram{\begin{tikzcd}[ampersand replacement=\&,row sep=small,column sep=small]
\& G/R \arrow[r] \& G/B\\
\deleted{T}\arrow[r] \& \overset{\cup}{T/R} \&\\
\& G\arrow[r,dash] \& G/B\arrow[uu]\\
\& \overset{\cup}{T} \&
\end{tikzcd}}

\diagram{\begin{tikzcd}[ampersand replacement=\&,row sep=small]
\operatorname{Pic}(G/B) \arrow[d,equal] \& \deleted{\unsure{\operatorname{Pic}}}D(T/R) \arrow[l,hook'] \arrow[d]\\
\operatorname{Pic}(G/B) \& \deleted{\unsure{\operatorname{Pic}}}D(T)\arrow[l]
\end{tikzcd}}

\diagram{\begin{tikzcd}[ampersand replacement=\&,row sep=tiny]
G/R=\operatorname{GP}(n)\arrow[r] \& \unsure{\mathcal D_n|\mathcal D(n)|\mathrm{Drap}|\text{Dém.}}\\
\deleted{\unsure{\operatorname{Gl}|\operatorname{G}}} \&\\
G=\operatorname{Gl}(n)\arrow[r] \& \mathcal D_n\\
\overset{\cup}{T(n)} \&
\end{tikzcd}}

Si $G$ est \deleted{\unsure{réductive|réductive et}} \inserted{semi-simple}, $\operatorname{Pic}(G)=0$ sss $G$ est simplement connexe.

\uline{Proposition} $G$ \deleted{\unsure{réductive sur|réductif un}} groupe alg affine sur corps $k$. Alors $\exists$ entier $n>0$ [égal au degré d'une extension convenable $k'$ de $k$] tel que $\unsure{n\cdot|n}\operatorname{Pic}(G)=0$. \deleted{Si $G_{\mathrm{réd}}$ \deleted{est} est \deleted{résoluble}\inserted{unipotent}, \unsure{alors} alors on peut prendre pour $n$ une puissance de la \unsure{caract.}}

On \unsure{a} les précisions suivantes
% archive sheet number: 13
